% !PRINTED 196
\weberpage{196}

Über den Grad in bezug auf $\varkappa^2$ läßt sich noch schließen, daß die
Koeffizienten von $x^{2\nu}$ den Grad $\nu$ in bezug auf $\varkappa^2$ nicht
übersteigen. Denn ist diese Regel richtig für $n$ und $n+1$, so folgt ihre
Richtigkeit für $2n$ und $2n+1$, und für $n=1$, $n=2$ trifft sie zu.

Aus den Grundgleichungen zwischen den drei elliptischen Funktionen:
\[
1 = \mathrm{cn}^2 v + \mathrm{sn}^2 v = \mathrm{dn}^2 v + \varkappa^2\,\mathrm{sn}^2 x % [sic]
\]
ergeben sich noch die Relationen:
\weq{15}{
\begin{array}{l}
D^2 = B^2 + x^2y^2z^2A^2 = C^2 + \varkappa^2x^2y^2z^2A^2,\quad n\text{ gerade,}\\
D^2 = y^2B^2 + x^2A^2 = z^2C^2 + \varkappa^2x^2A^2,\quad n\text{ ungerade,}
\end{array}}
mit deren Hilfe man nach (14) $B_{2n}$ und $C_{2n}$ durch $A_n$ und $D_n$
allein so ausdrücken kann:
% [sic] The source prints (14), but the displayed recurrences for B_{2n}
% and C_{2n} are in (13); (14) concerns the odd indices 2n+1.
\weq{16}{
\begin{array}{r@{\;}c@{\;}l@{\qquad}l}
B_{2n} &=& D_n^4 - 2x^2y^2z^2A_n^2D_n^2 + \varkappa^2x^4y^4z^4A_n^4 & \\
C_{2n} &=& D_n^4 - 2\varkappa^2x^2y^2z^2A_n^2D_n^2 + \varkappa^2x^4y^4z^4A_n^4 & \left.\right\}\ n\text{ gerade,}\\[2mm]
B_{2n} &=& D_n^4 - 2x^2A_n^2D_n^2 + \varkappa^2x^4A_n^4 & \\
C_{2n} &=& D_n^4 - 2\varkappa^2x^2A_n^2D_n^2 + \varkappa^2x^4A_n^4 & \left.\right\}\ n\text{ ungerade.}
\end{array}}

\begin{center}
\textbf{\S\ 58. \gesperrt{Multiplikation der Funktion} $\wp(u)$.}
\end{center}

Eine sehr elegante Form erhält die Multiplikationstheorie für die
\gesperrt{Weierstrasssche} Funktion $\wp(u)$ mit den beiden Perioden $\omega_1,
\omega_2$.

Betrachten wir die doppelt periodische Funktion
\weq{1}{\wp(nu) - \wp(u),}
worin $n$ eine beliebige positive ganze Zahl sein mag. Diese Funktion wird
unendlich für $u=0$, und zwar so, daß
\weq{2}{\wp(nu) - \wp(u) + \frac{n^2-1}{n^2}\,\frac{1}{u^2}}
für $u=0$ endlich bleibt. Außerdem wird sie aber unendlich für alle Werte
$u^*$ von $u$, die der Bedingung
\[
n\,u^* \equiv 0 \pmod{\omega_1,\omega_2}
\]
genügen, also von der Form sind:
\weq{3}{u^* = \frac{\nu_1\omega_1 + \nu_2\omega_2}{n},}
worin $\nu_1, \nu_2$ beliebige ganze Zahlen bedeuten. Wir erhalten alle in
(3) enthaltenen inkongruenten Werte, wenn wir $\nu_1$ und $\nu_2$ je ein
vollständiges Restsystem nach dem Modul $n$ durchlaufen lassen.
